%==================================================================
%  Quarto -> Beamer 用プリアンブル v2（16:9。v1 はリポジトリ直下の _extensions/koyama/）
%  koyama2.lua がこのファイルをヘッダに入れる。その直前に、YAML から作った設定マクロが入る:
%    \koyamaNoSectionOutline … section-outline: false のとき定義される
%    \koyamaBibFile          … 文献ファイル（bib-file:、既定 ref.bib）
%    \koyamaEvent / \koyamaVenue … 表紙右上の学会名・会場（event: / venue:）
%    \koyamaAuthorCount / \koyamaAuthors … 表紙の著者（author: の name / affiliation / presenter）
%  palette: で差し替えた色の \definecolor もここより前に入る。
%
%  v2 のデザイン
%    - 白地。帯やナビバーは無し。スライドのタイトルの左にセクション番号の四角（影つき）
%    - 右上にセクション名（グレー）とページ番号。セクション名の下に、そのセクションの中での
%      位置を ○● で示す
%    - 文字は全部太字（Noto Sans / Noto Sans JP の Bold。\textbf は ExtraBold）
%    - 表紙・目次・右上の表示は headline の中で、ページの左上を原点にして描く
%==================================================================

%========================
%  パッケージ
%========================
\usepackage{xcolor}
\usepackage{colortbl}      % \cellcolor（表セルの塗り）用
\usepackage{graphicx}
\usepackage{tabularx}
\usepackage{multirow}
\usepackage{booktabs}
\usepackage{amsmath, amssymb}
\usepackage{pgfplots}
\pgfplotsset{compat=1.18}
\usepackage{tikz}
\usetikzlibrary{shapes.geometric, arrows.meta, calc, positioning, patterns.meta}
\usepackage{tcolorbox}
\tcbuselibrary{skins}
\usepackage{caption}
\captionsetup[figure]{labelformat=empty, skip=2pt, font=small}
\usepackage{appendixnumberbeamer}

%========================
%  色の定義
%  役割: primary（ティール）=構造＋文中強調 / accent（コーラル赤）=文中強調・対比の一方
%        highlight（橙）=表・数値の注目点、3つ目の区分
%  ※ \providecolor なので、YAML の palette: で先に定義された色があればそちらが使われる。
%========================
\providecolor{textmain}{RGB}{34,34,34}        % 本文
\providecolor{textmuted}{RGB}{108,117,128}    % 出典 (Author, Year)・Note:
\providecolor{textlight}{RGB}{150,150,150}    % 右上のセクション名・ページ番号・○●
\providecolor{textfaint}{RGB}{190,190,190}    % 目次の、今のセクション以外
\providecolor{background}{RGB}{255,255,255}
\providecolor{primary}{RGB}{69,176,176}       % 構造の主役（四角・箇条書きの記号・目次の線）
\providecolor{primarylight}{RGB}{217,239,239} % 変数定義ボックスの地
\providecolor{accent}{RGB}{237,80,87}         % 文中強調・対比の一方・枠（Research Gap など）
\providecolor{accentlight}{RGB}{250,205,207}
\providecolor{highlight}{RGB}{232,155,0}      % 表の注目セル・3つ目の区分
% 主色から作る中間の色
\colorlet{primarymid}{primary!45!white}       % 四角の影
\colorlet{primarypale}{primary!30!white}      % 目次の、今のセクション以外の四角
\colorlet{primarypaler}{primary!13!white}     % その影

%========================
%  テーマ（metropolis を土台に、ヘッダー・表紙・目次は下で作り直す）
%========================
\usetheme{metropolis}

\setbeamercolor{normal text}{fg=textmain, bg=background}
\setbeamercolor{background canvas}{bg=background}
\setbeamercolor{frametitle}{fg=textmain, bg=background}
\setbeamercolor{alerted text}{fg=accent}
\setbeamercolor{block title}{bg=primary, fg=white}
\setbeamercolor{block body}{bg=primarylight}
\setbeamercolor{block title alerted}{bg=accent, fg=white}
\setbeamercolor{block body alerted}{bg=accentlight}
\setbeamercolor{itemize item}{fg=primary}
\setbeamercolor{itemize subitem}{fg=primary}
\setbeamercolor{itemize subsubitem}{fg=primary}
\setbeamercolor{enumerate item}{fg=primary}
\setbeamercolor{enumerate subitem}{fg=primary}
\setbeamercolor{button}{bg=primary, fg=white}

\setbeamertemplate{navigation symbols}{}
\setbeamertemplate{footline}{}
\setbeamersize{text margin left=11.3mm, text margin right=11.3mm}
% 本文は上寄せ（タイトルのすぐ下から書く）
\makeatletter\beamer@centeredfalse\makeatother

% \appendix が part page（"Appendix" ボックス）を挿入しないよう無効化
\AtBeginPart{}
% allowframebreaks のタイトルに付く続きの番号（i, ii …）を付けない
\setbeamertemplate{frametitle continuation}{}

%========================
%  フォント（v1 と同じ。ただし欧文の本文は Regular ではなく Medium にして、和文の Medium とそろえる）
%  Latin=Noto Sans（通常=Medium／太字=Bold。等幅は Noto Sans Mono）
%  日本語=Noto Sans JP（通常=Medium／太字=SemiBold）
%  v2 では、スライドのタイトル・表紙・目次・右上の表示・囲みの見出しと中身・箇条書きの見出しを
%  太字にする（本文は通常の太さ）。
%========================
\usepackage[sfdefault, medium]{noto}
\usepackage[match]{luatexja-fontspec}
\setmainjfont{Noto Sans JP}[UprightFont={* Medium},BoldFont={* SemiBold}]
\setsansjfont{Noto Sans JP}[UprightFont={* Medium},BoldFont={* SemiBold}]
% 行送りを広めに（11pt に対して約 16.5pt）
\renewcommand{\baselinestretch}{1.2}

\setbeamerfont{frametitle}{size=\fontsize{13}{16}\selectfont, series=\bfseries}
\setbeamerfont{itemize/enumerate body}{size=\normalsize}
\setbeamerfont{itemize/enumerate subbody}{size=\normalsize}
\setbeamerfont{itemize/enumerate subsubbody}{size=\normalsize}

%========================
%  箇条書き
%  1段目＝ティールの小さい ●、2・3段目＝ティールの短い横線 -
%  2段目の記号は1段目の文字の頭にそろえる
%========================
\setbeamertemplate{itemize item}{\tikz[baseline={(0,-0.36em)}]\fill[primary] (0,0) circle (0.6mm);}
\setbeamertemplate{itemize subitem}{\textcolor{primary}{\rule[0.3em]{1.3mm}{0.42mm}}}
\setbeamertemplate{itemize subsubitem}{\textcolor{primary}{\rule[0.3em]{1.3mm}{0.42mm}}}
\setbeamertemplate{enumerate item}{\textcolor{primary}{\insertenumlabel.}}
\setbeamertemplate{enumerate subitem}{\textcolor{primary}{\insertenumlabel.\insertsubenumlabel}}
\setlength\leftmargini{7.4mm}
\setlength\leftmarginii{2.6mm}
\setlength\leftmarginiii{3.4mm}
\setlength\labelsep{1.2mm}
\makeatletter
\def\@listi{\leftmargin\leftmargini \labelsep 1.2mm
  \topsep 3\p@ \@plus2\p@ \@minus2.5\p@ \parsep 0\p@ \itemsep 2.5\p@ \@plus1\p@ \@minus1\p@}
\let\@listI\@listi
\def\@listii{\leftmargin\leftmarginii \labelsep 1.4mm
  \topsep 1\p@ \@plus1\p@ \@minus1\p@ \parsep 0\p@ \itemsep 1.5\p@ \@plus1\p@}
\def\@listiii{\leftmargin\leftmarginiii \labelsep 1.4mm
  \topsep 1\p@ \@plus1\p@ \@minus1\p@ \parsep 0\p@ \itemsep 1.5\p@ \@plus1\p@}
\makeatother

%========================
%  ヘッダー・表紙・目次
%  すべて headline の中で、ページの左上を原点に描く（headline の高さは 0）。
%  セクション内の枚数（○● の数）と目次のセクション名は、前回の実行で aux に書いたものを読む。
%  変わったときは "Rerun to get cross-references right" を出し、Quarto にもう一度 lualatex を回させる。
%========================
\makeatletter
\providecommand{\koyamaEvent}{}
\providecommand{\koyamaVenue}{}
\providecommand{\koyamaAuthorCount}{0}
\providecommand{\koyamaAuthors}{}
\providecommand{\koyamaOutlineName}{目次}

\newif\ifkoyama@titlepage
\newif\ifkoyama@outline
\newif\ifkoyama@changed
\newif\ifkoyama@appendix
\newif\ifkoyama@nonumber
\newif\ifkoyama@nosection   % 右上のセクション名と ○● を出さない（References）
\def\koyama@uncounted{-1}
% \appendix の時点のセクション番号。これより後のセクションが付録（A1, A2 …）
\def\koyama@appbase{9999}
% スライドのタイトルに入れると、四角に番号を書かない（koyama2.lua が References のスライドに入れる）
%  右上のセクション名と ○● も出さず、そのセクションの枚数にも数えない
\def\koyamaNoSectionNumber{\global\koyama@nonumbertrue\global\koyama@nosectiontrue}
\newcount\koyama@secframe   % 今のセクションの何枚目か
\newcount\koyama@maxsec
\newcount\koyama@i
\def\koyama@lastframe{-1}
\def\koyama@lastsec{-1}
\def\koyama@nsections{0}

% aux から読む（前回の実行の結果）
\def\koyama@secframes#1#2{\expandafter\xdef\csname koyama@nsec@#1\endcsname{#2}}
\long\def\koyama@secname#1#2{\expandafter\gdef\csname koyama@sname@#1\endcsname{#2}}
\def\koyama@secmax#1{\xdef\koyama@nsections{#1}}

\def\koyama@twodigit#1{\ifnum#1<10 0\fi\number#1}
% 四角に書くセクション番号：本編は 01, 02 …、付録は A1, A2 …
\def\koyama@seclabel#1{%
  \ifnum#1>\koyama@appbase A\number\numexpr#1-\koyama@appbase\relax\else\koyama@twodigit{#1}\fi}

% 新しいフレームに入ったら、セクション内の番号を数える（目次・表紙は数えない）
\def\koyama@countframe{%
  \ifnum\c@framenumber=\koyama@lastframe\relax\else
    \xdef\koyama@lastframe{\the\c@framenumber}%
    \ifnum\c@section=\koyama@lastsec\relax\else
      \global\koyama@secframe=0
      \xdef\koyama@lastsec{\the\c@section}%
      \ifnum\c@section>\koyama@maxsec \global\koyama@maxsec=\c@section\fi
    \fi
    \ifkoyama@outline\else\ifkoyama@titlepage\else
      \global\advance\koyama@secframe by 1
      \expandafter\xdef\csname koyama@cur@\the\c@section\endcsname{\the\koyama@secframe}%
    \fi\fi
  \fi}

\def\koyama@pagenumfont{\fontsize{12}{14}\selectfont\bfseries}
\def\koyama@squarefont{\fontsize{10}{12}\selectfont\bfseries}

% ページ番号（右上）
\def\koyama@pagenumber{%
  \node[anchor=east, inner sep=0pt, font=\koyama@pagenumfont, text=textlight]
    (koyama@pn) at (\paperwidth-5mm, -7.2mm) {\insertframenumber};}

% 通常のスライドの右上：セクション名（グレー）＋ その下に ○●（セクション内の位置）＋ ページ番号
\def\koyama@drawheader{%
  \begin{tikzpicture}[overlay]
    \koyama@pagenumber
    \ifkoyama@nosection\else\ifnum\c@section>0
      \path let \p1 = (koyama@pn.west) in
        coordinate (koyama@secr) at ({min(\paperwidth-12mm, \x1-2.5mm)}, -5.9mm);
      \node[anchor=east, inner sep=0pt, text=textlight,
            font=\fontsize{9}{11}\selectfont\bfseries\addfontfeature{LetterSpace=6}\ltjsetparameter{kanjiskip={0.12\zw}}]
        at (koyama@secr) {\secname};
      % ○● の数：前回の実行で数えた、このセクションの枚数（初回は今の位置まで）
      \expandafter\ifx\csname koyama@nsec@\the\c@section\endcsname\relax
        \edef\koyama@n{\the\koyama@secframe}%
      \else
        \edef\koyama@n{\csname koyama@nsec@\the\c@section\endcsname}%
        \ifnum\koyama@n<\koyama@secframe \edef\koyama@n{\the\koyama@secframe}\fi
      \fi
      \koyama@i=0
      \loop\ifnum\koyama@i<\koyama@n\relax
        \advance\koyama@i by 1
        \ifnum\koyama@i=\koyama@secframe
          \fill[textlight] ($(koyama@secr)+({-(\koyama@n-\the\koyama@i)*1.9mm-0.6mm}, -3.7mm)$) circle (0.6mm);
        \else
          \draw[textlight, line width=0.45pt]
            ($(koyama@secr)+({-(\koyama@n-\the\koyama@i)*1.9mm-0.6mm}, -3.7mm)$) circle (0.55mm);
        \fi
      \repeat
    \fi\fi
  \end{tikzpicture}}

% 目次（各セクションの冒頭に自動で入る）
%   左上にデッキの短いタイトル（shorttitle:、無ければ title）、左に「目次」とティールの縦線、
%   右に番号の四角とセクション名。今のセクションは濃い四角＋黒字、それ以外は淡い四角＋薄いグレー
%   並べるのは本編のセクションだけ（付録のセクションは載せない）
\def\koyama@drawoutline{%
  \begin{tikzpicture}[overlay]
    \node[anchor=west, inner sep=0pt, font=\fontsize{13}{16}\selectfont\bfseries, text=textmain]
      at (6.3mm, -7.2mm) {\insertshorttitle};
    \koyama@pagenumber
    \node[anchor=west, inner sep=0pt, font=\fontsize{13.5}{16}\selectfont\bfseries, text=textmain]
      at (11.3mm, -16.2mm) {\koyamaOutlineName};
    \draw[primary, line width=1.3pt] (26.6mm, -13mm) -- (26.6mm, -84.2mm);
    \ifnum\koyama@nsections>1
      \pgfmathsetlengthmacro\koyama@step{min(15.45mm, 61.8mm/(\koyama@nsections-1))}%
    \else
      \def\koyama@step{0pt}%
    \fi
    \pgfmathsetlengthmacro\koyama@sq{min(8.2mm, 0.8*\koyama@step+(\koyama@step==0)*8.2mm)}%
    \koyama@i=0
    \loop\ifnum\koyama@i<\koyama@nsections\relax
      \advance\koyama@i by 1
      \pgfmathsetlengthmacro\koyama@y{-17.1mm-(\the\koyama@i-1)*\koyama@step}%
      \ifnum\koyama@i=\c@section
        \def\koyama@sqc{primary}\def\koyama@shc{primarymid}\def\koyama@txc{textmain}%
      \else
        \def\koyama@sqc{primarypale}\def\koyama@shc{primarypaler}\def\koyama@txc{textfaint}%
      \fi
      \fill[\koyama@shc] (33.2mm+0.12*\koyama@sq, \koyama@y+0.5*\koyama@sq-0.14*\koyama@sq)
        rectangle ++(\koyama@sq, -\koyama@sq);
      \fill[\koyama@sqc] (33.2mm, \koyama@y+0.5*\koyama@sq) rectangle ++(\koyama@sq, -\koyama@sq);
      \node[inner sep=0pt, text=white, font=\koyama@squarefont]
        at (33.2mm+0.5*\koyama@sq, \koyama@y) {\koyama@twodigit{\koyama@i}};
      \node[anchor=west, inner sep=0pt, text=\koyama@txc, font=\fontsize{13}{16}\selectfont\bfseries]
        at (47.3mm, \koyama@y) {\csname koyama@sname@\the\koyama@i\endcsname};
    \repeat
  \end{tikzpicture}}

% 表紙
%   右上：学会名（ティール）＋ ティールと赤のグラデーションの帯 ＋ 日付@会場
%   中央：タイトル（2行まで）＋ 淡いティールの線 ＋ 著者（名前と所属を横に並べる。発表者に *）
%   最下部：本編のセクション名と、その下にスライドの数だけ ○ を並べた帯（全体の構成と枚数の目安）
\newcommand{\koyamaAuthor}[3]{}
\def\koyama@drawtitlepage{%
  \begin{tikzpicture}[overlay]
    \ifx\koyamaEvent\@empty
      \def\koyama@datey{-7.5mm}%
    \else
      \def\koyama@datey{-16.6mm}%
      \node[anchor=east, inner sep=0pt, text=primary, font=\fontsize{11.5}{14}\selectfont\bfseries]
        at (\paperwidth-8mm, -7.4mm) {\koyamaEvent};
      \shade[left color=white, right color=primary]
        (\paperwidth-65.3mm, -11.0mm) rectangle (\paperwidth-6.4mm, -11.8mm);
      \shade[left color=white, right color=accent]
        (\paperwidth-65.3mm, -11.8mm) rectangle (\paperwidth-6.4mm, -12.9mm);
    \fi
    \node[anchor=east, inner sep=0pt, text=textmain, font=\fontsize{11}{13}\selectfont\bfseries]
      at (\paperwidth-8mm, \koyama@datey)
      {\insertdate\ifx\koyamaVenue\@empty\else\ifx\beamer@date\@empty\else @\fi\koyamaVenue\fi};
    \pgfmathsetlengthmacro\koyama@tw{\paperwidth-24mm}%
    \node[anchor=south, inner sep=0pt, align=flush center, text width=\koyama@tw, text=textmain,
          font=\fontsize{14.5}{24}\selectfont\bfseries]
      at (0.5\paperwidth, -43.3mm)
      {\inserttitle\ifx\beamer@subtitle\@empty\else\\{\fontsize{12}{20}\selectfont\insertsubtitle}\fi};
    \draw[primary!65!white, line width=0.9pt] (13.8mm, -44.8mm) -- (\paperwidth-13.8mm, -44.8mm);
    \ifnum\koyamaAuthorCount>0
      \koyama@i=0
      \pgfmathsetlengthmacro\koyama@colw{(\paperwidth-27.6mm)/\koyamaAuthorCount}%
      \pgfmathsetlengthmacro\koyama@tw{\koyama@colw-2mm}%
      \def\koyamaAuthor##1##2##3{%
        \advance\koyama@i by 1
        \pgfmathsetlengthmacro\koyama@ax{13.8mm+(\the\koyama@i-0.5)*\koyama@colw}%
        \node[anchor=south, inner sep=0pt, align=flush center, text width=\koyama@tw,
              font=\fontsize{12}{14}\selectfont\bfseries, text=textmain]
          at (\koyama@ax, -58.0mm) {##1\ifnum##3=1 *\fi};
        \node[anchor=north, inner sep=0pt, align=flush center, text width=\koyama@tw,
              font=\fontsize{8.5}{10.5}\selectfont\bfseries, text=textmain]
          at (\koyama@ax, -58.9mm) {##2};}%
      \koyamaAuthors
    \fi
    \koyama@titleprogress
  \end{tikzpicture}}

% 表紙の最下部の帯：本編のセクションを等間隔に並べ、名前の下に、そのセクションのスライドの数だけ ○
% （v1 の上部のナビバーと同じ形。表紙なので今の位置の ● は無い）。数は前回のコンパイルで数えたもの
\def\koyama@titleprogress{%
  \ifnum\koyama@nsections>0
    \pgfmathsetlengthmacro\koyama@colw{(\paperwidth-12.6mm)/\koyama@nsections}%
    \koyama@i=0
    \loop\ifnum\koyama@i<\koyama@nsections\relax
      \advance\koyama@i by 1
      % 塗りの帯にはせず、セクションごとの欄の上端に細い淡いティールの線（タイトル下の線と同じ色）。
      % 線の両端を少し詰めて欄の間に隙間を作り、線の本数でセクションの数も見せる
      \draw[primary!65!white, line width=0.6pt]
        (6.3mm+\koyama@i*\koyama@colw-\koyama@colw+1.2mm, -82.2mm) --
        (6.3mm+\koyama@i*\koyama@colw-1.2mm, -82.2mm);
      \pgfmathsetlengthmacro\koyama@ax{6.3mm+(\the\koyama@i-1)*\koyama@colw+1.5mm}%
      \node[anchor=west, inner sep=0pt, text=textmain, font=\fontsize{8}{9.5}\selectfont\bfseries]
        at (\koyama@ax, -84.7mm) {\csname koyama@sname@\the\koyama@i\endcsname};
      \expandafter\ifx\csname koyama@nsec@\the\koyama@i\endcsname\relax
        \def\koyama@n{0}%
      \else
        \edef\koyama@n{\csname koyama@nsec@\the\koyama@i\endcsname}%
      \fi
      \ifnum\koyama@n>0
        \pgfmathsetlengthmacro\koyama@dstep{min(1.6mm, (\koyama@colw-3mm)/\koyama@n)}%
        \foreach \koyama@k in {1,...,\koyama@n}{%
          \draw[textmain, line width=0.4pt]
            (\koyama@ax+0.5mm+\koyama@k*\koyama@dstep-\koyama@dstep, -87.6mm) circle (0.5mm);
        }%
      \fi
    \repeat
  \fi}

\setbeamertemplate{headline}{%
  \koyama@countframe
  % References など右上を出さないスライドは、フレームの始めに数えた分を戻す（1枚につき1回）
  \ifkoyama@nosection
    \ifnum\c@framenumber=\koyama@uncounted\relax\else
      \xdef\koyama@uncounted{\the\c@framenumber}%
      \global\advance\koyama@secframe by -1
      \expandafter\xdef\csname koyama@cur@\the\c@section\endcsname{\the\koyama@secframe}%
    \fi
  \fi
  \vbox to 0pt{\hbox to 0pt{%
    \ifkoyama@titlepage
      \koyama@drawtitlepage
      \global\koyama@titlepagefalse
    \else\ifkoyama@outline
      \koyama@drawoutline
    \else
      \koyama@drawheader
    \fi\fi
    \hss}\vss}%
  \global\koyama@nosectionfalse}

% 表紙のフレーム：中身は headline で描くので、ここでは印を付けるだけ
\setbeamertemplate{title page}{\global\koyama@titlepagetrue\vfill\mbox{}}

% スライドのタイトル：左にセクション番号の四角（影つき）、続けて黒字のタイトル
% （beamer はタイトル欄を本文の左余白の位置から組むので、その分だけ左へ戻してページの左端を原点にする）
%   References のスライドは四角だけ（番号なし）。付録のセクションは A1, A2 …
\setbeamertemplate{frametitle}{%
  \nointerlineskip
  % タイトルの中の \koyamaNoSectionNumber を先に効かせるため、いったん組んで捨てる
  \setbox\z@\vbox{\hsize=10cm \insertframetitle}%
  \hbox to\textwidth{\hskip-\beamer@leftmargin
  \begin{tikzpicture}
    \useasboundingbox (0,0) rectangle (\paperwidth, -14mm);
    \ifnum\c@section>0
      \fill[primarymid] (7.3mm, -4.7mm) rectangle ++(8.2mm, -8.2mm);
      \fill[primary] (6.3mm, -3.5mm) rectangle ++(8.2mm, -8.2mm);
      \ifkoyama@nonumber\else
        \node[inner sep=0pt, text=white, font=\koyama@squarefont] at (10.4mm, -7.6mm)
          {\koyama@seclabel{\c@section}};
      \fi
      \def\koyama@tx{18.3mm}%
    \else
      \def\koyama@tx{6.3mm}%
    \fi
    \pgfmathsetlengthmacro\koyama@tw{\paperwidth-\koyama@tx-42mm}%
    \node[anchor=west, inner sep=0pt, align=flush left, text width=\koyama@tw,
          text=textmain, font=\usebeamerfont{frametitle}]
      at (\koyama@tx, -7.6mm) {\insertframetitle};
  \end{tikzpicture}\hss}%
  \global\koyama@nonumberfalse
  \par\vskip 1mm}

% 2つ目以降の各セクションの冒頭に目次を自動で入れる。YAML で section-outline: false のときは入れない。
% セクション名は aux に書いておき、次の実行で目次に並べる。
\def\koyama@recordsection{%
  \koyama@startappendix
  \immediate\write\@auxout{\string\koyama@secname{\the\c@section}{\unexpanded\expandafter{\secname}}}}
\ifdefined\koyamaNoSectionOutline
  \AtBeginSection[]{\koyama@recordsection}
\else
  \AtBeginSection[]{%
    \koyama@recordsection
    % 最初のセクションと付録のセクションには目次を入れない
    % （全体の構成は表紙の最下部の帯で見せる。目次は「次にどこへ行くか」を示すために2つ目から）
    \ifkoyama@appendix\else\ifnum\c@section>1
    \begingroup
    \global\koyama@outlinetrue
    \begin{frame}\mbox{}\end{frame}
    \global\koyama@outlinefalse
    \endgroup
    \fi\fi}
\fi

% \appendix は「次のセクションの始まりで付録に入る」予約にする。
% Quarto では \appendix が直前のスライド（References など）の中に入るので、その場で
% appendixnumberbeamer の \appendix を動かすと、そのスライドのページ番号が 0 になってしまう。
% 付録に入るときに、付録が始まる前のセクション番号を覚える（それより後が A1, A2 …）。
\newif\ifkoyama@appendixpending
\let\koyama@origappendix\appendix
\def\appendix{\global\koyama@appendixpendingtrue}
\def\koyama@startappendix{%
  \ifkoyama@appendixpending
    \global\koyama@appendixpendingfalse
    \global\koyama@appendixtrue
    \xdef\koyama@appbase{\the\numexpr\c@section-1\relax}%
    \koyama@origappendix
  \fi}

% 今回数えた枚数を aux に書く。前回と違えば、もう一度コンパイルさせる
\AtEndDocument{%
  \koyama@i=0
  \loop\ifnum\koyama@i<\numexpr\koyama@maxsec+1\relax
    \expandafter\ifx\csname koyama@cur@\the\koyama@i\endcsname\relax\else
      \immediate\write\@auxout{\string\koyama@secframes{\the\koyama@i}{\csname koyama@cur@\the\koyama@i\endcsname}}%
      \expandafter\ifx\csname koyama@nsec@\the\koyama@i\endcsname\relax
        \koyama@changedtrue
      \else
        \ifnum\csname koyama@nsec@\the\koyama@i\endcsname=\csname koyama@cur@\the\koyama@i\endcsname\relax
        \else\koyama@changedtrue\fi
      \fi
    \fi
    \advance\koyama@i by 1
  \repeat
  % 目次に並べる本編のセクションの数（付録があれば、付録が始まる前まで）
  \ifkoyama@appendix \koyama@maxsec=\koyama@appbase\relax\fi
  \immediate\write\@auxout{\string\koyama@secmax{\the\koyama@maxsec}}%
  \ifnum\koyama@nsections=\koyama@maxsec\relax\else\koyama@changedtrue\fi
  \ifkoyama@changed
    \@latex@warning@no@line{koyama2: section frames changed. Rerun to get cross-references right}%
  \fi}
\makeatother

%========================
%  tikz スタイル（決定木・DAG 用）
%========================
\tikzset{
  dag_node/.style={
    draw=textmain, ellipse,
    minimum width=2.0cm, minimum height=1.0cm,
    thick, align=center, fill=background, font=\small
  },
  box_node/.style={          % 通常のノード（ティールの枠）
    draw=primary, line width=1.1pt, fill=background,
    minimum width=2.6cm, minimum height=0.9cm, align=center, font=\small
  },
  key_node/.style={          % 強調したい結論ノード（赤の枠）
    draw=accent, line width=1.1pt, fill=background,
    minimum width=2.6cm, minimum height=0.9cm, align=center, font=\small
  },
  arrow/.style={
    ->, >=Stealth, line width=1.0pt, draw=textlight,
    shorten >=2pt, shorten <=2pt
  }
}

%========================
%  カスタム部品
%  囲み（defbox・tagbox・keytakeaway・bluebox）は本文の左右の余白より少し外まで広げる
%========================
% (1) 数式の変数定義ゾーン（淡いティールの地、枠なし）
%     Markdown で書くときは ::: defbox … :::（中の `- **見出し**` は ● なしの見出しになる）
\newtcolorbox{defbox}{
  enhanced, sharp corners, colback=primarylight, colframe=primarylight, boxrule=0pt,
  grow to left by=4.7mm, grow to right by=4.7mm,
  left=5mm, right=5mm, top=3mm, bottom=3mm, before skip=6pt, after skip=6pt
}

% (1b) 補足の囲み（primary の細い枠だけ）。ふだんは使わない
\newtcolorbox{bluebox}{
  colback=background, colframe=primary,
  boxrule=0.8pt, arc=1.5mm, left=8pt, right=8pt, top=4pt, bottom=4pt,
  before skip=4pt, after skip=4pt, fontupper=\small
}

% (1c) 見出しつきの囲み（Research Question / Research Gap など）
%      左上に丸いラベル、枠の中は中央ぞろえ。色は accent（既定）か primary
%      使い方: \begin{tagbox}{Research Gap} … \end{tagbox}
%              \begin{tagbox}[primary]{Research Question} … \end{tagbox}
\newtcolorbox{tagbox}[2][accent]{
  enhanced, colback=background, colframe=#1, boxrule=1.1pt, arc=1.8mm,
  grow to left by=4.7mm, grow to right by=4.7mm,
  halign=flush center, left=6mm, right=6mm, top=3.2mm, bottom=3.2mm,
  before skip=10pt, after skip=8pt, fontupper=\bfseries,
  title={#2}, coltitle=white, fonttitle=\fontsize{11.5}{13}\selectfont\bfseries,
  % ラベルは枠の上の線に半分かかる位置に。yshift*（上の余白をラベルの分だけ自動で足す）ではなく
  % yshift にして、枠の線から文字までの空きを上下で同じ（top=bottom）にする
  attach boxed title to top left={yshift=-\tcboxedtitleheight/2},
  boxed title style={colback=#1, colframe=#1, boxrule=0pt, arc=2.6mm,
    left=4mm, right=4mm, top=0.5mm, bottom=0.5mm}
}

% (2) 表の中の値の「highlight の枠」強調（文中には使わない）
\newtcbox{\ybox}{on line, colframe=highlight, colback=highlight!20,
  boxrule=1pt, arc=2pt, boxsep=1.5pt, left=2pt, right=2pt, top=1pt, bottom=1pt}

% (2b) 核心の一文（赤い枠・中央ぞろえ・少し大きく）。1デッキで数回まで
%      使い方: \begin{keytakeaway} … \end{keytakeaway}  /  \begin{keytakeaway}[primary] … \end{keytakeaway}
\newtcolorbox{keytakeaway}[1][accent]{
  enhanced, colback=background, colframe=#1, boxrule=1.1pt, arc=0.6mm,
  grow to left by=4.7mm, grow to right by=4.7mm,
  halign=flush center, valign=center, left=4mm, right=4mm, top=2.6mm, bottom=2.6mm,
  before skip=8pt, after skip=6pt, fontupper=\fontsize{12.5}{16}\selectfont\bfseries
}

% (2c) 斜線のマーカー（語句の下半分に斜線を引く）。文中のキーワード・リード文の要点に
%      使い方: \hatch{語句}（accent）/ \hatch[primary]{語句}。行をまたぐ語句には使えない
\newcommand{\hatch}[2][accent]{%
  \tikz[baseline=(koyamahatch.base)]{%
    \node[inner sep=0pt, outer sep=0pt, anchor=base west] (koyamahatch) at (0,0) {\phantom{#2}};
    \fill[overlay, pattern={Lines[angle=45, distance=1.7pt, line width=0.7pt]}, pattern color=#1!70!white]
      ([xshift=-0.8mm, yshift=-0.2em]koyamahatch.base west) rectangle
      ([xshift=0.8mm, yshift=0.36em]koyamahatch.base east);
    \node[inner sep=0pt, outer sep=0pt, anchor=base west] at (0,0) {#2};}}

% (2c') 中央ぞろえの1行（リード文・図の見出しなど）。使い方: \leadline{…}
\newcommand{\leadline}[1]{\par{\centering #1\par}\vspace{2pt}}

% (2d) 流れの下向き三角（▼）。前の内容から次の内容へ（中央に置く）
\newcommand{\flowdown}{%
  \par\vspace{2pt}{\centering\tikz\fill[textmain!90] (0,0) -- (5mm,0) -- (2.5mm,-1.9mm) -- cycle;\par}\vspace{2pt}}

% (2e) 値のカード（角の丸い色枠）。区分ごとの値を並べるときに
%      使い方: \cardbox[primary]{軽症\\3.84分}  \cardbox[highlight]{…}  \cardbox[accent]{…}
\newtcbox{\koyamacard}[1][primary]{on line, enhanced, colback=background, colframe=#1,
  boxrule=1.2pt, arc=1.5mm, left=5mm, right=5mm, top=1.5mm, bottom=1.5mm}
\newcommand{\cardbox}[2][primary]{\koyamacard[#1]{\begin{tabular}{@{}c@{}}#2\end{tabular}}}

% (2f) 箇条書きの太字の見出し（● なし・左端）。koyama2.lua が `- **見出し**`（子の箇条書きつき）を
%      これに置き換える。2つ目以降の見出しの前には \koyamaListGap で空きを入れる。
\newcommand{\koyamaListHead}[1]{\par\textbf{#1}\par\nobreak\vspace{-\parskip}}
\newcommand{\koyamaListGap}{\par\vspace{0.7em}}

% (3) Table / Figure ラベル（黒）
\newcommand{\tlabel}[1]{\textbf{#1}}

% (3b) 図表下の注釈・箇条書き項目の出典（グレー・小さめ）
\newcommand{\slidenote}[1]{{\footnotesize\color{textmuted}#1}}

% (3c) \slidenote よりさらに小さい注（表の下のSE・有意水準の注記など）
\newcommand{\tinynote}[1]{{\scriptsize\color{textmuted}#1}}

% (3d) 箇条書き項目の出典。(Author, Year) が途中で折り返さないよう \mbox で包む
\newcommand{\srcite}[1]{\slidenote{\mbox{\parencite{#1}}}}

% (4) 図のプレースホルダ（実画像が未配置でもレンダリング可能にする）
\newcommand{\figplaceholder}[2][7cm]{%
  \begin{tikzpicture}
    \node[draw=textlight, dashed, rounded corners=4pt, fill=black!2,
          minimum width=#1, minimum height=3.4cm,
          align=center, text=textlight, font=\small]
      {図をここに配置\\[2pt] \ttfamily #2};
  \end{tikzpicture}}

%========================
%  参考文献（biblatex / APA + biber）
%========================
\usepackage{csquotes}
\usepackage[style=apa, backend=biber]{biblatex}
\DeclareLanguageMapping{english}{english-apa}
\providecommand{\koyamaBibFile}{ref.bib}
\addbibresource{\koyamaBibFile}
\AtEveryBibitem{
  \clearfield{doi}\clearfield{url}\clearfield{isbn}\clearfield{issn}
  \clearfield{note}\clearfield{month}\clearfield{issue}\clearfield{number}
  \clearfield{pmid}\clearfield{pmcid}\clearfield{abstract}
}
% 文献リストは小さく・行間を詰めて、1枚になるべく多く載せる（デッキ側で \koyamaBibSize を先に定義すれば変えられる）
\providecommand{\koyamaBibSize}{\fontsize{7}{8.4}\selectfont}
\renewcommand*{\bibfont}{\koyamaBibSize}
\setbeamerfont{bibliography entry author}{size=\koyamaBibSize}
\setbeamerfont{bibliography entry title}{size=\koyamaBibSize, series=\mdseries}
\setbeamerfont{bibliography entry location}{size=\koyamaBibSize}
\setbeamerfont{bibliography entry note}{size=\koyamaBibSize}
\setbeamercolor{bibliography entry author}{fg=textmain}
\setbeamercolor{bibliography entry title}{fg=textmain}
\setbeamercolor{bibliography entry location}{fg=textmain}
\setbeamercolor{bibliography entry note}{fg=textmain}
\setlength{\bibitemsep}{1pt}
\setlength{\bibhang}{1.5em}
\setbeamertemplate{bibliography item}{}

%========================
%  ハイパーリンク
%========================
\hypersetup{colorlinks=true, urlcolor=primary, linkcolor=textmain, citecolor=textmuted}
